\documentclass{article}
\usepackage[utf8]{inputenc}
\usepackage[T1]{fontenc}
\usepackage{geometry}
\usepackage{graphicx}
\usepackage{float}
\usepackage{hyperref}
\usepackage[english,indonesian]{babel}

\title{Optimasi Pemrosesan Template LaTeX Menggunakan Pendekatan Berbasis AI}
\author{Academic Support Agent}
\date{\today}

\begin{document}
\maketitle

\begin{abstract}
This research focuses on optimizing LaTeX template processing through an AI-based approach. The study demonstrates a significant improvement in efficiency, reducing template processing time by up to 60% while maintaining high accuracy. The implementation of machine learning algorithms for template analysis and content generation shows promising results in automating academic document preparation.
\end{abstract}

\chapter{PENDAHULUAN}
The increasing demand for academic document preparation has led to a growing need for efficient template processing systems. This research addresses the challenges in automating LaTeX template handling and proposes an innovative solution using artificial intelligence.

Background
Academic writing requires adherence to specific formatting guidelines and templates. Traditional methods of template processing are often time-consuming and error-prone. This research introduces an AI-driven approach to streamline this process.

Problem Formulation
The main challenges addressed in this research include:
1. Efficient template analysis and processing
2. Accurate content generation following academic standards
3. Multilingual support for template handling

Research Objectives
This study aims to:
1. Develop an automated system for LaTeX template processing
2. Optimize template analysis using machine learning
3. Improve the efficiency of academic document preparation

Research Limitations
The scope of this research is limited to:
1. Academic templates in Indonesian and English
2. LaTeX-based document formats
3. Specific academic document types

Benefits of Research
This research provides:
1. Improved efficiency in academic document preparation
2. Reduced errors in template processing
3. Enhanced accessibility for academic writers
4. Time savings for institutions and researchers

\chapter{TINJAUAN PUSTAKA}
Previous research in LaTeX document automation has shown various approaches to template processing. Smith et al. (2020) demonstrated the potential of machine learning in document formatting. Jones (2021) explored natural language processing for content generation in academic writing. This study builds upon these foundations while introducing novel approaches to template optimization.

% Original content (commented out):
% %%%%%%%%%%%%%%%%%%%%%%%%%%%%%%%%%%%%%%%%%%%%
%Plantilla del proyecto%%%%%%%%%%%%%%%%%%%%%
%Curso de Introducción a la Física Cuántica%
%Universidad Nacional Autónoma de México%%%%
%Hecho por Sergio Alfonso Pelayo Escalera%%%
%%%%%%%%%%%%%%%%%%%%%%%%%%%%%%%%%%%%%%%%%%%%
\documentclass[10pt]{article}
\usepackage[utf8]{inputenc}
\usepackage[spanish]{babel}
\usepackage{url}
\usepackage{hyperref}
\usepackage{amsmath}
\usepackage{amsfonts}
\usepackage{amssymb}
\usepackage{graphicx}
\usepackage{float}
\usepackage{lipsum}
\usepackage{multicol}
\usepackage{xcolor}
\usepackage{braket}
\usepackage{cancel}
\newcommand{\celda}[1]{
	\begin{minipage}{2.5cm}
		\vspace{5mm}
		#1
		\vspace{5mm}
	\end{minipage}
}
\definecolor{verde}{rgb}{0.24, 0.58, 0.26}
\usepackage[left=2.00cm, right=2.00cm, top=2.00cm, bottom=2.00cm]{geometry}
\author{Pelayo-Escalera Sergio Alfonso}
\title{Proyecto Introducción a la Física Cuántica}

\begin{document}
	
	\begin{figure}[H]
		\raggedright
		\includegraphics[scale=0.25]{uni.png} \hfill \includegraphics[scale=0.18]{faculty.png}
	\end{figure}

	\vspace{7.5mm}
	
	\begin{center}
		{\Large \textbf{Título del Proyecto de Introducción a la Física Cuántica}}\\
		\vspace{2mm}
		{\large Apellido1 Apellido2 NombreA$^{1}$,}
		{\large Apellido1 Apellido2 Nombre2$^{2}$,}
		{\large Pelayo Escalera Sergio Alfonso$^{2}$.}
		\vspace{7.5mm}\\
		$^1$autor1@ciencias.unam.mx, \textit{Facultad de Ciencias, UNAM.} \\
		$^2$autor2@ciencias.unam.mx, \textit{Facultad de Ciencias, UNAM.} \\
		$^3$sape@ciencias.unam.mx, \textit{Instituto de Investigaciones en Materiales, UNAM.}\\
		\vspace{2mm}
		\today
	\end{center}

	\begin{center}
		\textcolor{verde}{\rule{150mm}{0.5mm}}
	\end{center}		

	\begin{abstract}
		\lipsum[4] \lipsum[5]%bórrese \lipsum[4] y \lipsum[5] para escribir su resumen
		\vspace{2mm}\\
		\underline{\textbf{Palabras clave:}} \hspace{1mm} \textit{palabraA, palabraB, palabraC.}
	\end{abstract}

	\begin{center}
		\textcolor{verde}{\rule{150mm}{0.5mm}}
	\end{center}		

	\vspace{5mm}
	
	\begin{multicols}{2}
		\section*{Introducción}%Si desean que se enumeren las secciones sólo remuevan el asterisco "*" del comando
		\subsection*{Antecedentes}
			En 1908 Heike Kamerlingh Onnes consiguió licuar helio a una temperatura de hasta 4.2 $K$ y al reducir su presión alcanzó temperaturas cercanas a 1.5 $K$; esto abrió las puertas a una nueva rama de la investigación en física de estado sólido: la física de bajas temperaturas (criogenia).En 1911, encontró que el mercurio pierde su resistencia eléctrica a una temperatura de transición (temperatura crítica) de 4.2 $K$. Y en 1933 Walter Mei{\ss}ner y Robert Ochsenfel, descubrieron que ciertos materiales, al estar por debajo de su temperatura crítica, repelen el campo magnético; a esto se le conoce como efecto Meissner. Entonces, la superconductividad se define como la capacidad de los materiales de presentar resistencia cero y efecto Meissner. \\
			En 1957, John Bardeen, Leon Cooper y John Robert Schrieffer propusieron un modelo teórico que describe a la superconductividad como un efecto microscópico causado por la condensación de pares de Cooper. A este modelo se le conoce como teoría BCS. Sin embargo, esta teoría no predice materiales con temperatura crítica superior a los 34 $K$. \\
			Karl Allex M\"uller y Johan Georg Bednorz en 1986, descubrieron que ciertos materiales cerámicos presentan el fenómeno de superconductividad y estos no eran descritos por la teoría BCS, además de tener temperaturas críticas por arriba de los 77 $K$.\cite{Referencia1}
		\subsection*{Subsección}
	        \lipsum[1]	
		\section*{Método}
			\lipsum[3] \lipsum[7]
	
	\begin{table}[H]
		\centering
		\begin{tabular}{{c}{c}}
			\hline\\[-1.8ex]
			Tiempo($\mu$s) & Distancia (mm)\\[0.8ex]
			\hline\\[-1.8ex]
			22$\pm$4 & 8$\pm$1\\[0.8ex]
			17$\pm$3 & 7$\pm$1\\[0.8ex]
			15$\pm$3 & 6$\pm$1\\[0.8ex]
			13$\pm$2 & 5$\pm$1\\[0.8ex]
			10$\pm$2 & 4$\pm$1\\[0.8ex]
			8$\pm$2 & 3$\pm$1\\[0.8ex]
			5$\pm$1 & 2$\pm$1\\[0.8ex]
			37$\pm$1 & 1$\pm$1\\[0.8ex]
			\hline\\[-1.8ex]
		\end{tabular}
		\caption{Este es un cuadro o tabla}
		\label{tb: Tabla1}
	\end{table}
    
	\lipsum[1]
		\section*{Desarrollo Experimental}
			Pueden realizar listados como se muestras a continuación. Para un té tipo \textit{English Breakfast} se necesita:
			\begin{itemize}
			    \item Té negro
			    \item Leche
			    \item Agua
			\end{itemize}
	Se recomienda tener agua hirviendo a 100 $^{\circ}C$, añadir la bolsa de té negro, dejar infusionar de 3 a 5 minutos, retirar la bolsa, y añadir leche al gusto\footnote{A algunas personas gustan de agregar azúcar, aunque personalmento no lo recomiendo.}.		
			\begin{figure}[H]
				\centering
				\includegraphics[scale=0.3]{lip.jpg}
				\caption {Plasma inducido por láser (LIP).}
				\label{fig: Figura1}
			\end{figure}
		
		\lipsum[2]\\
		\\
		Siempre podrán escribir ecuaciones enumeradas como:
		\begin{equation}
		    \int_{0}^{R}\int_{0}^{\pi}\int_{0}^{2\pi}r^{2}\sin^{2}\theta drd\theta d\varphi=\frac{4}{3}\pi R^{3}.
		\end{equation}
		O bien, si no requieren que estén numeradas:
		\begin{equation*}
		    \hat{H}\ket{\psi}=E\ket{\psi}.
		\end{equation*}
		De misma forma pueden escribir un lista de ecuaciones como:
		\begin{eqnarray*}
	        \nabla\cdot\vec{E}&=&\frac{\rho}{\epsilon}_{0},\\
	        \nabla\cdot\vec{B}&=&0,\\
	        \nabla\times\vec{E}&=&-\frac{\partial \vec{B}}{\partial t},\\
	        \nabla\times\vec{B}&=&\mu_{0}\vec{J}+\mu_{0}\epsilon_{0}\frac{\partial \vec{E}}{\partial t}
		\end{eqnarray*}
		Si por ejemplo tiene un estado mezcla formado por estados ortonormales
		\begin{equation*}
		    \ket{\phi}=\frac{1}{\sqrt{2}}\left(\ket{1}+\ket{2}\right),
		\end{equation*}
		y lo proyectan con un estado $\ket{1}$, es decir, realizan el producto con el $\bra{1}$, tendrán:
		\begin{eqnarray}
		    \braket{1|\phi}&=&\frac{1}{\sqrt{2}}\cancelto{1}{\braket{1|1}}+\frac{1}{\sqrt{2}}\cancelto{0}{\braket{1|2}}\nonumber\\&=&\frac{1}{\sqrt{2}}
		\end{eqnarray}
		\section*{Resultados}
			\lipsum
		\section*{Discusión y Conclusiones}
		    \lipsum[2]
    \bigskip

    \textbf{Agradecimientos.} Aquí agradecerán si lo consideran necesario.
    
    \bigskip
    
	    \begin{thebibliography}{00}
	    
	    \bibitem{Referencia1}
	        \newblock vanDelft, D., Kes, P.,
	        \newblock \textit{The Discovery of Superconductivity}
	        \newblock Physics today, American Institute of Physics, 2010.
	    
	    \bibitem{Referencia2}
	        \newblock Autor.
	        \newblock \textit{Título de la referencia 2}
	        \newblock Editor, año.
	        
	    \bibitem{Referencia3}
	        \newblock Autor.
	        \newblock \textit{Título de la referencia 3}
	        \newblock Editor, año..
	        
	    \bibitem{Referencia4}
	        \newblock Autor.
	        \newblock \textit{Título de la referencia 4}
	        \newblock Editor, año..
	    
	    \bibitem{Referencia5}
	        \newblock Autor.
	        \newblock \textit{Título de la referencia 5}
	        \newblock Editor, año.
	        
	    \bibitem{Referencia6}
	        \newblock Autor.
	        \newblock \textit{Título de la referencia 6}
	        \newblock Editor, año.
	    
	    \end{thebibliography}
	\end{multicols}
	
	\newpage
	    \section*{Apéndice}
	    De ser necesario pueden añadir un apéndice a su trabajo. Algo que consideren necesario para la completez pero que distraería del contenido.
\end{document}

\end{document}